% ID: SEL-D0002-OB-E03
% ESTADO: cotejada. Fuente y transcripción contrastadas de forma independiente; salida revisada.
% Fuente: Selectividad/sel_2019_matematicas_aplicadas/Titular Septiembre/Exámenes Mat Apl CCSS II - A y B.pdf
% SHA-256: 31fb46947e141f5c8da8f7b381e7d27b67fe5a9d1597120454ce66fce0313f1a
% Página: 2; opción B; ejercicio 3.
% Transcripción de primera pasada basada en la imagen de la página original.
% Se conserva la barra de P(A/B) y los complementos A^c y B^c de la fuente.
% Sin solución ni modificación de datos o condiciones.

De dos sucesos \(A\) y \(B\) de un mismo espacio muestral se sabe que:
\[
P(A\cap B)=0.2\qquad P(A\cup B)=0.4\qquad P(A/B)=0.8
\]
\begin{enumerate}[label=\alph*),leftmargin=2.4em]
  \item \textbf{(1.2 puntos)} Calcule \(P(B)\) y \(P(A)\).
  \item \textbf{(0.5 puntos)} \textquestiondown{}Son los sucesos \(A\) y \(B\) independientes? Razone la respuesta.
  \item \textbf{(0.8 puntos)} Calcule \(P(A^c\cup B^c)\).
\end{enumerate}
